\chapter{Proofs for Classical Cut Elimination}
\label{annex:p07-classical-cut-elimination}

\begin{annexitem}{annex:p07.classical.alpha-identity}{thm:p07-classical-alpha-identity}{Exact identity without an alpha-conversion rule}
\begin{proof}
We first prove the syntactic facts needed for the construction. Encode a
formula tree by retaining predicate, function and logical-constructor
labels. At a variable leaf, record its actual name if it is free; if it
is bound, record the position of its nearest governing binder in the
stack of surrounding binders. The nearest-binder convention includes
shadowed names: a quantifier inside \(QxC\) with its own binder \(x\)
governs its own occurrences. A generating fresh alpha-renaming changes
the outer binder's name and just its governed leaves, so it leaves this
tree unchanged. Congruence, symmetry and transitivity preserve equality
of the encodings as well.

Conversely, suppose the two encodings agree. Traverse their corresponding
binders from outside inward and rename each corresponding pair to the same
variable, choosing these variables pairwise distinct and absent from all
original names and free-variable names of the two finite formulas. Every
step is an allowed fresh alpha-generator: the chosen name is absent from
the matrix at that step, including the newly chosen earlier binder names.
Capture-avoiding substitution stops at an inner binder of the renamed
variable, so shadowed occurrences are unchanged. Both resulting literal
trees now have the same binder labels, bound leaves, free leaves and other
labels; hence they are identical formulas. The finite sequences of
generators establish alpha-equivalence of the originals. Thus equality of
these encodings characterizes precisely the stated alpha-equivalence.

In particular alpha-equivalent formulas have the same outer constructor,
and corresponding operands of a binary constructor are alpha-equivalent.
For \(QxC=_{\alpha}QzD\), choose \(y\) outside every variable in the
two formulas. Removing the outer binder and replacing its governed leaves
by the free name \(y\) gives the same encoding on both sides. Therefore
\(C[y/x]=_{\alpha}D[y/z]\), with the exact least-fresh substitutions
on the left and right.

For a finite simultaneous substitution \(\sigma\), first use the preceding
freshening construction to give every binder a name outside the domain
of \(\sigma\) and every variable in its finitely many range terms.
Substitution now replaces exactly the free variable leaves by their range
terms, once per leaf, without capture. It commutes with function symbols,
falsity and each binary logical constructor. At a quantifier it leaves the
fresh binder in place and continues through its matrix. These cases exhaust
the term and formula grammars. Thus this operation is determined by the
encoding, so substitution respects alpha-equivalence. The exact algorithm
may choose different fresh names, but only bound-name labels change;
its result has the same encoding. A bijection of variable names similarly
acts only on free-name labels and the names in substituted terms; bound
positions remain fixed. This proves the commutation of permutation and
substitution on alpha-classes.

The substitution interchange used below can also be checked leaf by leaf.
Suppose \(x\) is outside the domain of \(\sigma\) and outside all variables
in its explicitly assigned range terms. Then
\[
 (C[t/x])\sigma=_{\alpha}(C\sigma)[t\sigma/x].
\]
For the free leaf \(x\), both sides give \(t\sigma\). For a free
leaf \(v\ne x\), both give \(v\sigma\), since that term has no free
\(x\); an unchanged leaf \(v\) is distinct from \(x\) as well.
Function constructors reduce to their argument leaves. Give all formula
binders names fresh for the terms and substitutions in this identity.
Then falsity, each binary constructor and each fresh quantifier preserve
the identity by induction through their operands or matrix. Undoing the
freshening proves the displayed alpha-equality for arbitrary literal
representatives.

Finally, the particular identity needed for principal quantifier reduction
does not require \(x\) to be absent from the witness term \(t\).
Choose \(y\ne x\) outside all variables of \(C,t\). Then
\[
 (C[y/x])[t/y]=_{\alpha}C[t/x].
\]
A free \(x\) leaf goes first to \(y\), then to \(t\); every other
free leaf is fixed, since none was \(y\). No second substitution takes
place inside the inserted \(t\). The same constructor induction with
fresh binders proves the formula identity. Also, substituting for a name
not free in a formula fixes its encoding and hence its alpha-class.
All these facts concern syntax alone.

Now prove both orientations of the asserted identity simultaneously by
induction on formula rank. For atoms, alpha-equivalence is literal equality,
so use Id. Falsity is also literally itself and has its own identity.
Suppose \(A=A_1\land A_2\), \(B=B_1\land B_2\).
From the smaller identities \(A_i\Rightarrow B_i\), left-weaken by the
other \(A\)-operand and apply the sheet's two-operand \(\land L\).
This gives \(A\Rightarrow B_1\) and \(A\Rightarrow B_2\).
The shared-context \(\land R\) gives \(A\Rightarrow B\).

For \(A=A_1\lor A_2\), \(B=B_1\lor B_2\), right-weaken
\(A_i\Rightarrow B_i\) by the other \(B\)-operand and apply the
single, two-auxiliary \(\lor R\). The two results
\(A_1\Rightarrow B\), \(A_2\Rightarrow B\) combine by \(\lor L\).
For implications, use the smaller identities
\(B_1\Rightarrow A_1\) and \(A_2\Rightarrow B_2\).
Split \(\to L\) gives \(A_1\to A_2,B_1\Rightarrow B_2\), and
\(\to R\) gives \(A_1\to A_2\Rightarrow B_1\to B_2\).

For \(A=\forall xC\), \(B=\forall zD\), take fresh \(y\).
The smaller alpha-identity is
\(C[y/x]\Rightarrow D[y/z]\), using the exact specified substitution
values. Apply \(\forall L\) with term \(y\), followed by
\(\forall R\) with eigenvariable \(y\), to obtain \(A\Rightarrow B\).
The final conclusion has no free \(y\).
For \(A=\exists xC\), \(B=\exists zD\), use the same smaller identity,
then \(\exists R\) on the right and \(\exists L\) on the left, both
with \(y\). Its eigenvariable condition also holds. Interchanging
\(A,B\) in each construction proves the opposite orientation.
No Cut, Mix, semantic equivalence or contexted identity axiom was used.
\end{proof}
\end{annexitem}

\begin{annexitem}{annex:p07.classical.parameter-substitution}{thm:p07-classical-parameter-substitution}{Height-preserving operations on alpha-class trees}
\begin{proof}
Work only in the auxiliary presentation. A bijection of variable names,
applied to every formula, witness term and eigenvariable annotation of a
proof, preserves its tree. The syntactic facts established above make the
quantified instances commute with this operation on classes. A bijection
also preserves absence from a free-variable set. Id, \(\bot L\), all four
structural rules and all six propositional rules remain the same instances,
including their shared or split contexts. All four quantifier rules retain
their instance and, where present, their freshness condition. This proves
the permutation assertion by induction on the tree.

At an eigenvariable inference, transpose its parameter \(y\) with a
variable \(z\) absent from the entire premise proof and the additional
finite data. Apply that permutation to the whole premise proof. Since
\(y\) is not free in the inference's conclusion and \(z\) is fresh,
the concluding classes are unchanged. This freshening preserves auxiliary
height. There are always enough variables because every specified proof
and forbidden set is finite.

For term substitution, the following slightly stronger induction is useful:
any finite simultaneous substitution \(\sigma\) applied to the endsequent
of a cut-free auxiliary tree has a cut-free tree of the same height.
Terms in its range are inserted only for free variables. Id maps to
\([C\sigma]\Rightarrow[C\sigma]\); \(\bot L\) retains falsity.
WL, WR, CL and CR are reapplied to their transformed occurrences, even
if previously different formulas now have the same class. Each of
\(\land L,\land R,\lor L,\lor R,\to L,\to R\) commutes with
substitution in its operands; shared contexts remain shared and split
contexts remain concatenated.

For \(\forall L\), display the quantified class as \(\forall xC\)
with bound \(x\) fresh for the domain and range variables of \(\sigma\).
This changes no class or proof tree. The transformed auxiliary instance is
\[
 [C[t/x]\,\sigma]=[(C\sigma)[t\sigma/x]].
\]
Apply the induction hypothesis to the premise and reapply \(\forall L\)
with term \(t\sigma\). For \(\exists R\), exactly the same displayed
identity supplies its right-hand instance and the same witness substitution.
These rules have no eigenvariable restriction.

For \(\forall R\), choose such a fresh bound representative and freshen
its eigenvariable to \(z\) outside the domain and range variables of
\(\sigma\) and all other finite data. Apply induction to the permuted
premise, whose height remains strictly below the original inference.
Its active formula is now \([(C\sigma)[z/x]]\). The substituted
conclusion is \(z\)-free: its old free variables did not include \(z\)
and the substitution range cannot introduce it. Hence \(\forall R\)
is valid. For \(\exists L\), freshen its premise parameter in the same
way; the active formula on the left becomes \([(C\sigma)[z/x]]\),
and the same free-variable calculation validates \(\exists L\).
Each case reapplies one original inference above premises of the same
height, completing the induction. One-variable substitution is included.

Term substitution changes terms, not logical constructors. The
bound-occurrence characterization likewise leaves all logical constructors
unchanged under alpha-conversion. Thus formula rank is preserved.
This proof makes no claim about the height of exact representative
realizations of these auxiliary trees.
\end{proof}
\end{annexitem}

\begin{annexitem}{annex:p07.classical.alpha-realization}{thm:p07-classical-alpha-realization}{Returning to the registered rule sheet}
\begin{proof}
An exact derivation becomes an auxiliary one by taking classes at every
occurrence; each inference and its height are unchanged. For the other
direction induct on the height of a cut-free auxiliary tree, with the
prescribed representative of every end occurrence included in the data.

At class identity, the two representatives are alpha-equivalent but need
not be literally equal. Use the exact cut-free alpha-identity just proved.
At \(\bot L\), the principal class has only the literal representative
\(\bot\), and the prescribed endsequent is an exact initial instance.
At weakening, realize the other end occurrences in the premise, then add
the prescribed new formula by WL or WR. At contraction, prescribe two
identical copies of its retained end representative in the premise,
realize the premise, then apply CL or CR. Other context occurrences keep
their prescribed representatives and multiplicities.

For \(\land L\), prescribe the two exact conjuncts of the chosen
principal in its premise. For \(\land R\), prescribe one such right
operand in each premise and the same exact context representatives in
both. For \(\lor L\), prescribe one left disjunct in each shared-context
premise. For \(\lor R\), prescribe both exact right disjuncts in its one
premise. For \(\to R\), prescribe the exact antecedent operand on the
left and consequent operand on the right. For split \(\to L\),
distribute the prescribed end context tokens according to the auxiliary
split and prescribe its exact antecedent operand in the first succedent
and consequent operand in the second antecedent. Constructor compatibility
of alpha-equivalence ensures these are the required premise classes.
Apply the realization induction to each premise and reapply that exact
registered rule.

For \(\forall L\), retain its witness term \(t\). If the prescribed
principal is \(\forall xC\), prescribe the exact, least-fresh
\(C[t/x]\) as its active premise occurrence. It has the required class
by the syntactic substitution facts. Realize that premise and apply
\(\forall L\). For \(\exists R\), use the same construction on the
right with its prescribed \(\exists xC\).
For \(\forall R\), freshen the auxiliary eigenvariable to \(y\)
outside every variable in the prescribed endsequent. This does not
change the auxiliary height or conclusion. Prescribe the exact instance
\(C[y/x]\) and the given contexts in the shorter premise; realize it
and apply the exact \(\forall R\). Its freshness condition holds for
the actual conclusion. For \(\exists L\), prescribe \(C[y/x]\) on
the left after the same freshening, realize the shorter premise, and
apply the exact \(\exists L\).

These are every cut-free rule. Realization is finite, contains no Cut or
Mix and has exactly the prescribed endsequent. Its identity replacements
can increase physical height, which is why realization is not inserted
inside a same-rank height comparison in the elimination induction.
\end{proof}
\end{annexitem}

\begin{annexitem}{annex:p07.classical.bottom-erasure}{thm:p07-classical-bottom-erasure}{Selected falsities, including initial sequents}
\begin{proof}
Induct on the height of a cut-free proof together with its selected
succedent bottom occurrences. With none selected, retain the proof.
At Id, a nonempty selection is possible only for
\(\bot\Rightarrow\bot\); replace it by the initial
\(\bot\Rightarrow\). At \(\bot L\), the antecedent bottom remains,
so the erased conclusion is another instance of that initial schema.

If WR introduces a selected token, erase the selected old tokens in
its premise and omit this WR. If the new token is unselected, erase
the old selection and reapply WR. If CR's retained token is selected,
select both ancestors, erase them by induction and omit CR. For an
unselected CR, retain both contracted ancestors, erase the other selected
tokens and reapply CR. WL and CL affect only the antecedent; erase the
selected succedent tokens in their premise and reapply those rules.

No right logical rule introduces bottom itself. All selected tokens in
a logical conclusion are therefore context tokens. For the unary rules
\(\land L,\lor R,\to R\), preserve every active operand, erase the
selected context ancestors and reapply the rule. For shared \(\land R\)
and \(\lor L\), erase the same selected context in both premises and
reapply the shared rule. For split \(\to L\), erase only the selected
ancestors assigned to each premise; reapply the split rule with its active
operands retained. For \(\forall L\) and \(\exists R\), the active
instance and witness term remain unchanged. For \(\forall R\) and
\(\exists L\), the active instance remains and deleting context formulas
cannot introduce a free eigenvariable into the conclusion. Thus all four
quantifier rules can also be reapplied.

An active operand equal to bottom is unselected unless it is an ancestor
of the designated end token; logical operands of retained principals are
never erased merely by comparing their spelling. This proves the
construction, with no increase in auxiliary height. The same induction
works for exact syntax, because it changes neither formulas nor witnesses.
\end{proof}
\end{annexitem}

\begin{annexitem}{annex:p07.classical.marked-mix}{thm:p07-classical-marked-mix}{The complete contraction-sensitive reduction}
\begin{proof}
Pass the two exact cut-free input trees to the auxiliary presentation.
All heights in this proof refer to auxiliary trees; realization occurs
only at the end. Keep the selected occurrence tokens as part of the data.
Write the inputs and required output as
\[
 D:\Gamma\Rightarrow\Delta,A^m,\qquad
 E:A^n,\Pi\Rightarrow\Sigma,
 \qquad \Gamma,\Pi\Rightarrow\Delta,\Sigma. \tag{M}
\]
If \(m=0\), use D and weaken with \(\Pi,\Sigma\). If \(n=0\), use
E and weaken with \(\Gamma,\Delta\). These include every zero-boundary
case. Otherwise induct on
\((\operatorname{rk}_{\mathrm{fm}}(A),\operatorname{ht}(D)+
\operatorname{ht}(E))\), universally over contexts, marked subsets and
multiplicities. The hypothesis places no upper bound on the height of
its constructed output.

\emph{Finite structural bookkeeping.}
Weakening adds any requested finite multiset by one step per occurrence.
If \(S\) contains \(r\) copies of a formula B, contracting \(j\) copies
of S into one means reducing \(jr\) occurrences of B to \(r\), not to
one. Context tokens and active tokens are tracked separately until the
active ones have been eliminated. Thus overlapping contexts and equal
active formulas do not authorize excess contraction.

\emph{All initial cases.}
If D is Id, its selected succedent forces D to be
\(A\Rightarrow A\), with \(m=1\), \(\Gamma=A\), \(\Delta\) empty.
Contract the n selected antecedent copies in E to one. The result
\(A,\Pi\Rightarrow\Sigma\) is (M), with every unselected copy retained.
If E is Id, it is \(A\Rightarrow A\), with \(n=1\), \(\Pi\) empty
and \(\Sigma=A\). Contract D's m selected copies to one to obtain (M).
If D is \(\bot L\), its left bottom lies in \(\Gamma\) and (M) is
itself initial. If E is \(\bot L\) with a retained antecedent bottom,
that token lies in \(\Pi\), so (M) is again initial. The remaining
initial possibility has E's every antecedent bottom selected. Then
\(A=\bot\); erase D's selected right bottoms by the preceding lemma,
then weaken with \(\Pi,\Sigma\). These are every initial pair.

\emph{All structural cases.}
If D ends with WR introducing a selected A, apply induction to its
immediate premise and the original E with multiplicity \(m-1,n\), and
omit that WR. If E ends with WL introducing a selected A, use the
original D and its immediate premise with \(m,n-1\), omitting that WL.
Zero multiplicity is allowed in both calls.
If D ends with CR forming a selected A, both ancestors are selected;
the call uses its immediate premise with \(m+1,n\), and no CR is
reapplied. If E ends with CL forming a selected A, use \(m,n+1\)
on its immediate premise and omit CL. In every case the selected input
height decreases, regardless of the increased multiplicity.

For an unselected WR or CR in D, keep its new occurrence or both
contracted ancestors unselected, transport the other marks, apply
induction to the premise and reapply that WR or CR. For an unselected
WL or CL in E, remove the other selection by the recursive call, then
restore the new unselected occurrence by WL, or retain both unselected
contraction ancestors and reapply CL to them.
WL and CL in D act on its antecedent, never its selected succedent:
apply induction to the premise and reapply the rule. WR and CR in E
act on its succedent, never its selected antecedent: do the same.
This exhausts the four structural rules on each side, including an
unselected contraction whose formula is also A. Exchange requires no
inference because contexts are multisets.

\emph{Nonselected logical principals.}
Suppose the last rule on one side does not introduce a selected token.
On the D side apply the induction to each immediate D-premise and the
original full E. On the E side apply it to the original full D and each
immediate E-premise. Marks lift through shared contexts to both premises
and through split contexts only to their assigned premise. A branch with
no marked ancestor uses the zero-boundary construction. Every call has
the same formula rank and a strictly smaller original input-height sum.

The effect is to insert the fixed contexts \((\Pi;\Sigma)\) in each
transformed D-premise, or \((\Gamma;\Delta)\) in each transformed
E-premise, while removing its marked context tokens. Denote the inserted
pair in this paragraph by \((L;R)\). Reconstruct each logical rule as
follows; its active operand tokens are unselected and remain present.
\begin{description}
\item[\(\land L\).] The premise retains its two displayed left operands;
form their conjunction once, with L and R as added contexts.
\item[\(\land R\).] Both premises acquire the same L and R and retain
their respective right operand. Their enlarged contexts agree, so apply
the shared-context conjunction rule.
\item[\(\lor L\).] Both premises retain their distinct left disjunct and
the same enlarged contexts; apply the shared-context disjunction rule.
\item[\(\lor R\).] Its one premise retains both right disjuncts. Form
their disjunction with the added contexts unchanged.
\item[\(\to R\).] Its one premise retains the displayed left antecedent
operand and right consequent operand. Apply implication-right to them.
\item[\(\to L\).] Apply the calls separately to its two split premises.
Reapplying split implication-left yields two copies of L and R. Contract
those two copied context multisets to one, preserving their own internal
multiplicities.
\item[\(\forall L\).] Retain the displayed left instance \(B[t/x]\) and
reapply universal-left with the same witness term t.
\item[\(\forall R\).] Before the recursive call, freshen its parameter
\(y\) for the entire opposite proof. Retain the right instance \(B[y/x]\) and
reapply universal-right; the inserted contexts are y-free.
\item[\(\exists L\).] Before the call, freshen its parameter y for the
opposite proof. Retain the left instance \(B[y/x]\) and reapply
existential-left with its valid freshness condition.
\item[\(\exists R\).] Retain the right instance \(B[t/x]\) and reapply
existential-right with the same term t.
\end{description}

For the split calculation on D, write its premises as
\[
 \Gamma_1\Rightarrow\Delta_1,A^{m_1},B,
 \qquad C,\Gamma_2\Rightarrow\Delta_2,A^{m_2},
 \quad m_1+m_2=m.
\]
The recursive results are
\(\Gamma_1,\Pi\Rightarrow\Delta_1,\Sigma,B\) and
\(C,\Gamma_2,\Pi\Rightarrow\Delta_2,\Sigma\).
Split implication-left concludes
\[
 B\to C,\Gamma_1,\Gamma_2,\Pi,\Pi
 \Rightarrow\Delta_1,\Delta_2,\Sigma,\Sigma.
\]
Contract the two added \(\Pi\) and \(\Sigma\) copies. On E, its split
contexts partition the n marked antecedent occurrences; the same
calculation duplicates \(\Gamma\) and \(\Delta\), which are then
contracted. Shared \(\land R\) and \(\lor L\) do not need this step:
their enlarged contexts occur in both premises but only once in the
conclusion of the shared rule.

To make the quantifier sides explicit, \(\forall L\) on D and
\(\exists R\) on E are always nonselected-principal cases, because
their principals lie on the wrong side for the selection. The same is
true of \(\exists L\) on D and \(\forall R\) on E, but these require
the stated eigenvariable freshening. The other four possibilities are
\(\forall R\) on D, \(\exists R\) on D, \(\forall L\) on E and
\(\exists L\) on E: the preceding reconstruction applies precisely
when their principal token is unselected. Thus all eight side-specific
quantifier commutations are included. Freshening preserves the original
auxiliary premise height before the same-rank call. It occurs before
introducing the other proof's contexts, not after an invalid inference.

\emph{Preparation of every principal pair.}
The remaining case has one selected principal A at the last logical
rule of each proof. Put \(S=\Gamma,\Pi\) and \(T=\Delta,\Sigma\).
There may still be selected A occurrences in their side contexts.
For each immediate premise \(D_i\), remove its residual marked A
occurrences by a call against the \emph{original full E}. For each
immediate premise \(E_j\), remove its residual marked occurrences by
a call against the \emph{original full D}. The inequalities are
\[
 \operatorname{ht}(D_i)+\operatorname{ht}(E)
 <\operatorname{ht}(D)+\operatorname{ht}(E),
 \qquad
 \operatorname{ht}(D)+\operatorname{ht}(E_j)
 <\operatorname{ht}(D)+\operatorname{ht}(E).
\]
If the residual class is empty, use the zero case. Only \emph{after}
these calls return may one add missing split-context portions by
weakening to the common S,T. These new proof heights are unrestricted;
they will be inputs only to smaller-formula-rank calls. Every lower-rank
call below selects exactly its displayed active tokens, leaving equal
formulas in S,T unselected.

\emph{Conjunction \(A=B\land C\).}
The last rules are \(\land R\) on D and \(\land L\) on E. Residual
cleaning produces
\[
 S\Rightarrow T,B,\qquad S\Rightarrow T,C,\qquad B,C,S\Rightarrow T.
\]
Use the smaller-rank Mix hypothesis on the displayed B in the first
and third proofs. Its result is \(C,S,S\Rightarrow T,T\).
Mix the displayed C with the second proof, obtaining
\(S,S,S\Rightarrow T,T,T\). Contract the three context copies.
Both B and C have smaller rank. If B=C, there are nevertheless two
different active tokens in the third premise, and they are consumed
one at a time.

\emph{Disjunction \(A=B\lor C\).}
The sheet uses one two-auxiliary \(\lor R\) premise and two shared
\(\lor L\) premises. Cleaning produces
\[
 S\Rightarrow T,B,C,\qquad B,S\Rightarrow T,\qquad C,S\Rightarrow T.
\]
Smaller-rank Mix on B gives \(S,S\Rightarrow T,T,C\); Mix on C with
the third proof gives three copies of S,T, which are contracted. No
single-injection disjunction-right rule was assumed.

\emph{Implication \(A=B\to C\).}
D's immediate premise is
\(B,\Gamma\Rightarrow\Delta,A^{m-1},C\).
E's split premises can be written
\[
 A^{n_1},\Pi_1\Rightarrow\Sigma_1,B,\qquad
 C,A^{n_2},\Pi_2\Rightarrow\Sigma_2,\qquad n_1+n_2=n-1,
\]
with \(\Pi=\Pi_1,\Pi_2\) and \(\Sigma=\Sigma_1,\Sigma_2\).
The prescribed same-rank calls, followed by the missing-context
weakenings, produce
\[
 B,S\Rightarrow T,C,\qquad S\Rightarrow T,B,\qquad C,S\Rightarrow T.
\]
Mix the second and first on B, giving
\(S,S\Rightarrow T,T,C\). Mix with the third on C, then contract
the three context copies. These two calls have smaller rank and do
not require a height bound on the cleaned proofs.

\emph{Universal \(A=\forall xB\).}
Freshen D's last eigenvariable y away from the whole E proof, its
witness t, every variable of B, the binder x and the retained contexts.
The last premises are
\[
 D_0:\Gamma\Rightarrow\Delta,A^{m-1},B[y/x],\qquad
 E_0:B[t/x],A^{n-1},\Pi\Rightarrow\Sigma.
\]
The calls \((D_0,E;m-1,n)\) and \((D,E_0;m,n-1)\) on A give
\[
 F:S\Rightarrow T,B[y/x],\qquad G:B[t/x],S\Rightarrow T.
\]
The parameter y is free in neither S nor T nor \(\forall xB\).
Freshen F's internal eigenvariables for t, and apply the independent
substitution lemma with \([t/y]\). Its endsequent is
\(S\Rightarrow T,B[t/x]\): the retained contexts are fixed on
classes and the capture-avoiding composition identity gives the active
instance. Mix this with G on that one instance and contract the two
copies of S,T. Its rank is
\(\operatorname{rk}_{\mathrm{fm}}(B[t/x])=
\operatorname{rk}_{\mathrm{fm}}(B)<
\operatorname{rk}_{\mathrm{fm}}(\forall xB)\).

\emph{Existential \(A=\exists xB\).}
Freshen E's last eigenvariable y away from D, t, every variable of B
and the binder x. Its premises are
\[
 D_0:\Gamma\Rightarrow\Delta,A^{m-1},B[t/x],\qquad
 E_0:B[y/x],A^{n-1},\Pi\Rightarrow\Sigma.
\]
The two smaller-height calls on A give
\(F:S\Rightarrow T,B[t/x]\) and \(G:B[y/x],S\Rightarrow T\).
Now substitute \([t/y]\) in G, after freshening its internal
eigenvariables. The result is \(B[t/x],S\Rightarrow T\), with
unchanged retained contexts. A smaller-rank Mix with F followed by
contraction gives (M). This is the dual placement of the substitution,
not an assertion that eigenvariables may occur in the added contexts.

There is no logical principal pair for an atom or for bottom, since
no right logical rule introduces either. The initial bottom case was
already settled independently. The five cases exhaust all constructors;
defined negation is the implication-to-bottom case.

Every recursive call has either the same rank and a smaller
\emph{original input} height sum, or a strictly smaller formula rank
with arbitrary finite input heights. Structural padding and contraction
after a call introduce no further same-rank call on that larger output.
Well-founded recursion on pairs of natural numbers therefore constructs
a finite auxiliary cut-free proof of (M). Finally realize it with the
exact prescribed context formulas and multiplicities. This proves the
marked Mix assertion on the registered sheet.
\end{proof}
\end{annexitem}

\begin{annexitem}{annex:p07.classical.cut-elimination}{thm:p07-classical-cut-elimination}{Eliminating all cuts in a finite tree}
\begin{proof}
First replace every formula occurrence by its alpha class. Normalize the
resulting finite proof by structural recursion on its original tree.
Initial inferences are retained. At any ordinary inference, recursively
normalize its immediate premises and reapply the original rule. Their
endsequents have not changed, so its context and eigenvariable conditions
still hold. At Cut, normalize its two premises and apply the marked Mix
construction with one selected copy on each side. At a registered Mix,
do the same with its displayed positive multiplicities. In either case
the returned proof is cut-free and has the original inference's conclusion.

The outer recursion passes only to immediate subtrees. Each inner Mix
construction terminates by its own proved lexicographic recursion.
Consequently the result is a finite auxiliary proof with neither Cut
nor Mix. Realize its end classes using the exact original endsequent
representatives. The realization theorem gives the claimed exact proof.
Its possibly greater physical height is irrelevant to the completed
auxiliary induction. No semantic completeness or separate cut theorem
has been used.
\end{proof}
\end{annexitem}

\begin{annexitem}{annex:p07.classical.analytic}{thm:p07-classical-analytic}{Subformula tracing and empty-sequent underivability}
\begin{proof}
In a cut-free proof, trace any formula occurrence downward toward the
endsequent. A context occurrence persists. A contracted occurrence joins
the retained occurrence of the same formula. An active propositional
operand becomes an immediate subformula of its principal formula.
Weakening creates an occurrence but does not remove any premise
occurrence. The shared rules supply a continuation from either premise,
and the split implication rule supplies the continuation in its assigned
context or principal operand. Thus every premise occurrence has a
continuation; no Cut or Mix is available to erase the trace.

At a quantifier rule, an instance \(B[t/x]\) or \(B[y/x]\) traces to the matrix
\(B\) of its principal quantifier. This is the required instance relation,
with alpha-renaming for bound names. Composing these finitely many trace
steps places every first-order occurrence in the closure of the end
formulas under taking subformulas, term instances and alpha-renaming.
No quantifier constructor present in a traced formula can disappear on
the path: a matrix instance has the same logical constructors as its
matrix, and the principal adds its displayed quantifier. Therefore a
cut-free proof of a propositional endsequent has no quantified formula
occurrences. Its traces use literal propositional subformulas only.
Cut elimination supplies such a cut-free proof from any derivation,
giving analytic existence after expanding defined negation.

Finally, the empty sequent is not an instance of Id or \(\bot L\).
Every other cut-free registered rule has a displayed formula in its
conclusion: a weakening formula, a retained contraction formula, or a
logical principal. Thus it cannot be the last inference of a cut-free
proof of the empty sequent. Cut elimination excludes an arbitrary
derivation of that sequent as well.
\end{proof}
\end{annexitem}
